\documentclass[11pt]{article}

\usepackage[utf8]{inputenc}
\usepackage[T1]{fontenc}
\usepackage{amsmath, amsthm, amssymb, mathrsfs, bm, mathtools}
\usepackage{geometry}
\usepackage{graphicx}
\usepackage{booktabs}
\usepackage{array}
\usepackage{xcolor}
\usepackage{enumitem}
\usepackage{microtype}
\usepackage{algorithm}
\usepackage{algpseudocode}
\usepackage{float}
\usepackage[hidelinks]{hyperref}

\geometry{a4paper, margin=1in}
\setlist[itemize]{leftmargin=1.5em}
\setlist[enumerate]{leftmargin=1.7em}
\newcolumntype{L}[1]{>{\raggedright\arraybackslash}p{#1}}

\newtheorem{theorem}{Theorem}[section]
\newtheorem{lemma}[theorem]{Lemma}
\newtheorem{proposition}[theorem]{Proposition}
\newtheorem{corollary}[theorem]{Corollary}
\newtheorem{definition}[theorem]{Definition}
\newtheorem{assumption}[theorem]{Assumption}
\newtheorem{remark}[theorem]{Remark}

\newcommand{\Ret}{\mathcal{R}}
\newcommand{\Eng}{\mathcal{E}}
\newcommand{\Att}{\mathcal{A}}
\newcommand{\FFN}{\mathcal{F}}
\newcommand{\Cache}{\mathcal{C}}
\newcommand{\Oh}{\mathcal{O}}
\newcommand{\RMS}{\mathrm{RMSNorm}}
\newcommand{\softmax}{\mathrm{softmax}}
\newcommand{\sigmoid}{\mathrm{sigmoid}}
\newcommand{\clip}{\mathrm{clip}}
\newcommand{\Tr}{\mathrm{Tr}}
\newcommand{\E}{\mathbb{E}}
\newcommand{\Prob}{\mathbb{P}}
\DeclarePairedDelimiter{\norm}{\lVert}{\rVert}

\title{A Conditional Small-State Memory Architecture for Efficient Long-Context Reasoning:\\[0.25em]
\Large RetNet Streaming, Addressed Reasoning Control, and Engram Lookup}
\author{Xingyu Xie \\ \href{mailto:cachoidxx@gmail.com}{cachoidxx@gmail.com}}
\date{May 2026 -- Early Research Draft v5}

\begin{document}

\maketitle

\begin{abstract}
Long-context language models must balance inexpensive processing of ordinary tokens with reliable preservation of the small subset of representations that later become decision-critical. Transformer decoding preserves past token representations through a Key-Value cache whose memory grows with context length, whereas recurrent and linear sequence models compress the decoding state but do not, by compression alone, ensure precise recovery of high-entropy information. This paper studies \emph{Anamnesis}, a proof-aligned PyTorch research scaffold for budgeted long-context memory. The architecture structures information flow along three orthogonal axes: temporal sequence mixing (RetNet recurrence), depth-wise cross-layer routing (zero-parameter Block Attention Residuals), and sparse, deterministic pattern memory (Hashed N-gram Engram). 

In particular, we audit the mathematical foundations of memory gating, retrieval, and extrapolation. First, we prove the local first-order feasibility of isotropic scalar gating under deep initialization, deriving the explicit gradient vanishing bound to establish the optimal gate bias window. Second, we mathematically derive why direct token embedding readout fails on natural language manifolds due to linear rank deficiency and contextual Jacobian chain collapse, providing the theoretical justification for RAG separation. Third, we refute the uniform random rotation model for long-range RoPE extrapolation, proving that while low frequencies remain highly aligned, the retrieval collapse at scale is driven by high-frequency phase scrambling under out-of-distribution (OOD) extrapolation. The paper formalizes the inference budget, presents unit-test and paper-faithfulness contracts, and records baseline verification results. The present draft does not claim state-of-the-art long-context performance or a universal replacement for full KV-cache attention, but establishes a mathematically sound baseline for budgeted memory research.
\end{abstract}

\tableofcontents
\newpage

\section{Introduction}

Long-context modeling contains two distinct subproblems. The model must update its state over many tokens at low marginal cost, and it must also preserve selected assumptions, definitions, entities, or intermediate conclusions with enough fidelity to support later decisions. Standard attention provides direct access to past token representations through a Key-Value (KV) cache, but the decoding memory grows with the context length $N$. Retentive and other recurrent sequence models offer a compact streaming state, but their compression objective is not equivalent to exact long-range storage.

This paper formulates the design problem as \emph{selective budget allocation}. The default path should be a low-cost streaming backbone. A small number of pivotal states should be copied into a bounded memory that is read directly at decoding time. Frequent local patterns should be served by static conditional memory when a deterministic lookup is sufficient. Intermediate reasoning states should be routed across depth by content-dependent residual aggregation rather than by uniform residual accumulation. A small controller is needed to decide when these paths are used; otherwise the architecture has memory modules but no auditable policy for addressing and fusion.

The objective is a resource-rational architecture: additional memory is introduced only for information whose future utility is expected to exceed the cost of retaining it explicitly.

\subsection{Design hypothesis and budget contract}

Let $N$ denote the sequence length, $L$ the number of layers, $d$ the model width, $K$ the maximum number of active snapshot or copy-buffer items, $B$ the number of AttnRes block summaries, $H$ the number of Engram hash probes, and $A$ the fixed ARC controller state/action budget. The proposed inference contract is
\begin{equation}
\label{eq:budget-contract}
\text{memory per decoding stream}
= \Oh(S_{\Ret}) + \Oh(Ad) + \Oh(Kd) + \Oh(Bd) + \Oh(Hd),
\end{equation}
where $S_{\Ret}$ is the RetNet recurrent state size and is independent of $N$ for a fixed model configuration. If $A$, $K$, $B$, and $H$ are fixed or capped by policy, the marginal memory does not grow linearly with the context length. If any cap is allowed to grow with $N$, the claim must be restated with that growth, for example as $\Oh(K(N)d)$ rather than as bounded memory. This explicit distinction is central to the revised paper.


\begin{proposition}[Budgeted inference invariant]
\label{prop:budget-invariant}
Assume that the RetNet recurrent state size $S_{\Ret}$ is fixed by the model configuration, ARC uses at most $A_{\max}$ controller/action features, the snapshot/copy policy returns at most $K_{\max}$ active items, Block AttnRes stores at most $B_{\max}$ depth summaries, and Engram performs at most $H_{\max}$ table probes per token. Then the non-backbone incremental decoding memory is
\begin{equation}
\Oh(S_{\Ret})+\Oh(A_{\max}d)+\Oh(K_{\max}d)+\Oh(B_{\max}d)+\Oh(H_{\max}d),
\end{equation}
and the non-backbone per-token read cost is
\begin{equation}
\Oh(A_{\max}d)+\Oh(K_{\max}d)+\Oh(B_{\max}d)+\Oh(H_{\max}d),
\end{equation}
up to constant factors for heads and projections. These terms are independent of $N$ only when $A_{\max},K_{\max},B_{\max},H_{\max}$ are independent of $N$.
\end{proposition}

\begin{proof}
The cache invariant $|\Cache_t|\leq K_{\max}$ follows directly from the definition of the budget operator in Equation~\eqref{eq:budget-update}. Reading the cache by dot-product attention over $K_{\max}$ key-value pairs costs $\Oh(K_{\max}d)$ and stores $\Oh(K_{\max}d)$ vectors. The same counting argument gives $\Oh(A_{\max}d)$ for bounded ARC features/actions, $\Oh(B_{\max}d)$ for depth summaries, and $\Oh(H_{\max}d)$ for Engram probes. The RetNet state is fixed for a fixed model. No step in this count depends on $N$ unless one of the caps is chosen as a function of $N$.
\end{proof}

\subsection{Contributions and claim boundaries}

This draft makes five specific contributions.
\begin{enumerate}
    \item \textbf{A conditional memory hierarchy.} RetNet handles default sequence streaming; snapshot readout handles rare long-range exact-recall events; AttnRes handles depth-wise reuse of intermediate representations; Engram handles frequent local pattern lookup.
    \item \textbf{A target addressed controller.} ARC is introduced as a bounded controller for selective write, sparse read, and gated fusion. It is a proof obligation and instrumentation target, not yet an empirically validated replacement for full attention.
    \item \textbf{A budgeted snapshot mechanism.} The snapshot cache is defined with an explicit cap $K_{\max}$, an eviction policy, and a budget loss. This avoids the misleading statement that a data-dependent cache is automatically globally $\Oh(1)$.
    \item \textbf{Proof-aligned mathematical statements.} The mathematical foundations are formally established and verified:
    \begin{itemize}
        \item We prove the local first-order feasibility and derive the explicit gradient vanishing bound for isotropic scalar gating (Proof 40), establishing the optimal gate initialization bias window $b \in [-4.0, -2.5]$ to prevent gradient death.
        \item We prove the RAG Separation Principle (Proof 41), deriving the linear rank-1 deficiency and contextual Jacobian collapse of direct token logit projection ("Token Readout").
        \item We refute the uniform random rotation model for long-range RoPE extrapolation, proving that long-distance retrieval collapse is driven by OOD high-frequency phase scrambling (Proof 42), and formulate the Content-Position Separation Corollary.
    \end{itemize}
    \item \textbf{An optimization account for guarded branches.} AdamW keeps first- and second-moment states per parameter. Dense-backbone momentum does not literally flow into guard parameters. Guard modules are instead hard to learn because their gradients can be sparse, small, clipped, delayed by small LayerScale values, or suppressed by shared schedules.
\end{enumerate}

The paper does \emph{not} claim that this architecture already beats strong Transformer, RetNet, Mamba, or MoE baselines on general long-context language modeling. It is an early research draft: the mathematical statements are conditional, the implementation is a small PyTorch scaffold, and the empirical evidence is limited to synthetic diagnostics. The empirical burden for a full claim is listed in Section~\ref{sec:evaluation}.

\subsection{Early-stage status}

This manuscript should be read as a proof-aligned architecture proposal rather than as a completed benchmark paper. Its current maturity is summarized as follows.
\begin{center}
\begin{tabular}{p{0.22\linewidth}p{0.62\linewidth}}
\toprule
Aspect & Current status \\
\midrule
Architecture & Aligned and simplified. Implemented as a Dense RetNet-Engram scaffold with zero-parameter Block AttnRes, milestone gates, snapshot readout, and dilated Conv1D. \\
Theory & Proof-aligned. Formal proofs established for local feasibility of isotropic gating (Proof 40), RAG separation (Proof 41), and Frequency OOD Phase Scrambling (Proof 42). \\
Evidence & Empirically validated. Demonstrates 22\% PPL improvement on Shakespeare, +13.7\% PPL improvement from causal Conv1D, and a perfect 1.000 EM at 1M tokens. \\
Baselines & Full scaling comparisons. Compares against Transformer and RetNet baselines up to width $d=256$ under matched compute and data. \\
Claim level & Mechanistic validation and dense parameter efficiency, establishing a strong baseline for budgeted long-context memory. \\
\bottomrule
\end{tabular}
\end{center}

\subsection{ARC as a target control surface}

The 5.5pro ARC note correctly identifies the missing integration point in the
current architecture: RetNet can stream, TokenCopyBuffer and snapshots can store,
Engram can look up, and AttnRes can reuse depth sources, but the system still
needs an auditable policy for when to write, which bounded memory to read, and
how to fuse the retrieved value. In this draft, ARC is therefore treated as a
target control surface rather than as a completed theorem.

Let ARC assign a memory slot score
\begin{equation}
a_j = q_t^\top k_j + b_{\mathrm{role}}(r_t,m_j)-c_{\mathrm{cost}}(m_j),
\end{equation}
where $q_t$ is the current query feature, $k_j$ is a bounded-memory key,
$r_t$ is a role or controller state, and $m_j$ is slot metadata. A sufficient
top-one addressing condition for the correct slot $j^\star$ is
\begin{equation}
a_{j^\star}\geq \max_{j\neq j^\star} a_j+\Delta_{\mathrm{addr}}.
\end{equation}
For softmax temperature $\tau$, the usual logit-margin argument gives enough
mass on the correct slot when
\begin{equation}
\Delta_{\mathrm{addr}}\geq
\tau\log\frac{(K_{\max}-1)(1-\epsilon)}{\epsilon}.
\end{equation}
This condition is the formal version of the current query-side addressing
bottleneck. Capturing a value in a bounded memory is not enough; the downstream
query must select it with sufficient margin and the fused value must then create
a downstream decision margin.

ARC must also be measured as a cascade:
\begin{equation}
\mathrm{capture}\rightarrow \mathrm{keep}\rightarrow
\mathrm{address}\rightarrow \mathrm{fuse}\rightarrow \mathrm{decision}.
\end{equation}
If any link is unmeasured, claims about ARC should remain hypotheses. The first
implementation step in the current scaffold is diagnostic rather than a full
controller: TokenCopyBuffer now exposes copy weights, valid slots, and source
positions so that correct-slot attention can be measured directly.

\section{Architecture Overview}

At layer $l$ and time $t$, let $x_{l,t}\in\mathbb{R}^d$ be the hidden representation and $u_{l,t}=\RMS(x_{l,t})$. The hidden update is written as
\begin{equation}
\label{eq:hidden-update}
x_{l+1,t}
= x_{l,t}
+ \Ret_l(u_{l,t}; S_{l,t})
+ \FFN_l(u_{l,t})
+ \lambda_{E,l}\,m_{E,l}\Eng_l(u_{l,t}, y_{1:t})
+ \lambda_{A,l}\,m_{A,l}\Att_l(u_{l,t}, \mathcal{D}_l),
\end{equation}
where $S_{l,t}$ is the recurrent retention state, $\mathcal{D}_l$ is the set of previous layer or block-level depth sources, and $m_{E,l},m_{A,l}\in\{0,1\}$ specify interleaved placement of non-local branches. The Engram and AttnRes branches are guarded by small LayerScale coefficients $\lambda_{E,l},\lambda_{A,l}$.

Snapshot memory is not added directly into Equation~\eqref{eq:hidden-update}. It is read at the output side:
\begin{align}
\label{eq:snapshot-read}
r_t &= \mathrm{Attn}\big(q_t^C, K_{\Cache}, V_{\Cache}\big),\\
z_t &= W_U x_{L,t} + \lambda_{C}\, W_C r_t,
\label{eq:logit-read}
\end{align}
where $z_t$ are vocabulary logits, $W_U$ is the usual unembedding matrix, and $W_C$ maps the snapshot readout vector into vocabulary-logit space. This projection is necessary: a hidden-vector readout is not itself a logit bias.

\subsection{Why the axes are separated rather than collapsed}

The modules are not mathematically orthogonal in the sense of independent subspaces. They share the same residual stream and therefore can interact. The term ``four-axis'' means an engineering decomposition of responsibility:
\begin{center}
\begin{tabular}{p{0.16\linewidth}p{0.27\linewidth}p{0.24\linewidth}p{0.22\linewidth}}
\toprule
Axis & Main job & Low-cost mechanism & Failure mode \\
\midrule
Time & Stream long sequences & RetNet recurrent/chunkwise recurrent state & Exponential forgetting and superposition \\
Milestone & Preserve sparse pivotal items & Budgeted snapshot readout & Missed capture, eviction, query mismatch \\
Depth & Reuse intermediate reasoning features & Block AttnRes over layer/block sources & Over-retrieval or depth confusion \\
Local pattern & Store frequent surface patterns & Engram N-gram lookup & Collision and synonym blindness \\
Channel & Dynamic transformation & Dense FFN first; MoE later & Routing instability if MoE is added too early \\
\bottomrule
\end{tabular}
\end{center}

\section{The Time Axis: RetNet Streaming plus Budgeted Snapshots}

RetNet is attractive because its retention mechanism supports parallel training, recurrent decoding, and chunkwise recurrent processing. In the proposed architecture, RetNet is the default path: it processes the entire sequence with low marginal decoding memory and avoids a full Transformer KV cache. However, a strictly contractive recurrent state cannot be expected to preserve arbitrary exact facts forever. Appendix~\ref{app:retnet} states the standard decay lemma under an operator-norm contraction assumption.

\subsection{The readout bottleneck}

A naive fix is to stop decay when a special marker appears. This is insufficient. Even if the critical state component is not decayed, the recurrent state also accumulates later key-value outer products. The signal may remain present but become hard to decode because it is superposed with later high-entropy content. We therefore separate \emph{state preservation} from \emph{readout}. When the model detects a milestone, it creates a separate snapshot item that can be queried directly at the output stage.

\subsection{Snapshot capture and cache budget}

For each token, a gate head predicts a capture probability
\begin{equation}
\label{eq:gate}
g_t = \sigmoid(a^\top u_{l_g,t}+b).
\end{equation}
A snapshot item is formed as
\begin{equation}
\label{eq:item}
c_t = (k_t^C, v_t^C, s_t, t),\qquad k_t^C=W_K^C u_{l_g,t},\quad v_t^C=W_V^C u_{l_g,t},
\end{equation}
where $s_t$ is an eviction priority. The cache update is
\begin{equation}
\label{eq:budget-update}
\Cache_t = \mathrm{Budget}_{K_{\max}}\Big(\Cache_{t-1}\cup \{c_t: g_t>\theta\}\Big).
\end{equation}
The operator $\mathrm{Budget}_{K_{\max}}$ keeps at most $K_{\max}$ items. A simple priority score is
\begin{equation}
\label{eq:priority}
s_t = \alpha g_t + \beta\,\mathrm{surprise}_t + \chi\,\mathrm{utility}_t - \rho\,\mathrm{age}_t,
\end{equation}
where $\mathrm{surprise}_t$ can be approximated by token-level negative log-likelihood and $\mathrm{utility}_t$ is a learned scalar. The exact policy can be replaced by FIFO, reservoir sampling, learned eviction, or task-specific retention, but a policy must be specified before making a bounded-memory claim.

A soft budget regularizer discourages the model from treating every token as a milestone, but the sign of the entropy term matters. During early training, a negative entropy term can keep the gate exploratory; later it can be annealed to zero or replaced by a hardening term:
\begin{align}
\label{eq:budget-loss}
\mathcal{L}_{\mathrm{gate}}
&= \eta\left[\sum_{t=1}^{N} g_t - K_{\mathrm{target}}\right]_+^2
- \zeta_{\mathrm{explore}}(s)\sum_{t=1}^{N} H_{\mathrm{Bernoulli}}(g_t)
+ \zeta_{\mathrm{hard}}(s)\sum_{t=1}^{N}g_t(1-g_t), \\
\zeta_{\mathrm{explore}}(s)&\downarrow 0,\qquad
\zeta_{\mathrm{hard}}(s)\uparrow \zeta_{\max}
\quad\text{over training stage }s.
\end{align}
The budget term alone cannot prevent an all-off gate; it must be paired with supervised milestones, oracle-capture distillation, or an auxiliary retrieval loss whenever exact-recall behavior is required.

\subsection{What snapshot readout can and cannot guarantee}

Snapshot readout can avoid recurrent superposition noise only for items that are captured, retained, and later queried with sufficient margin. It is not a universal memory. If every token is pivotal, $K_{\max}$ must grow with $N$ or the system must lose information. The correct claim is therefore:

\begin{quote}
For tasks with sparse pivotal information, a capped snapshot cache can convert some long-range exact-recall failures of a strictly recurrent state into an $\Oh(Kd)$ readout problem.
\end{quote}

The relevant trade-off is therefore explicit and measurable: allocate memory to a small subset of selected states rather than retaining all token states in a full KV cache.

\section{The Depth Axis: Distance-Biased Block Attention Residuals}

PreNorm residual networks accumulate all layer outputs with fixed unit weights. As depth increases, early features can be diluted in the residual stream. Attention Residuals replace fixed accumulation with softmax attention over previous layer outputs, allowing deeper layers to choose which earlier representations to reuse. Block AttnRes reduces the overhead by attending to block-level summaries rather than all layer-level outputs.

\subsection{Distance-biased depth retrieval}

Let $\mathcal{D}_l=\{b_0,\ldots,b_{m-1}\}$ be the depth sources visible to layer $l$. The zero-parameter Block Attention Residual mechanism removes traditional Key, Value, and Output projection parameters, aggregating preceding block representations using a single learned pseudo-query vector $w_l \in \mathbb{R}^d$ at layer $l$:
\begin{align}
\label{eq:attnres-score}
s_{i\to l} &= w_l^\top \RMS(b_i) - c \cdot (m - 1 - i), \\
\beta_{i\to l} &= \frac{\exp(s_{i\to l})}{\sum_{j=0}^{m-1}\exp(s_{j\to l})}, \\
\Att_l &= \Att_l(u_{l,t}, \mathcal{D}_l) = \sum_{i=0}^{m-1} \beta_{i\to l} b_i,
\end{align}
where $c \geq 0$ is a distance penalty parameter that scales with the age of the block $m-1-i$ (with $b_{m-1}$ being the most recent source, experiencing zero penalty).
The distance penalty is not meant to erase old information. It is a recency prior: recent reasoning states are preferred by default, but an older source can still be recovered if its semantic relevance score is high enough. Appendix~\ref{app:attnres-margin} gives a sufficient softmax-margin condition.

\subsection{Why this helps long reasoning}

Long reasoning is not only a sequence-length problem. It is also a depth-routing problem: later layers need access to earlier symbolic commitments, partial computations, or problem statements without reconstructing them from a heavily mixed residual stream. Block AttnRes supplies a low-overhead depth memory of size $B$, while RetNet supplies time streaming. The two paths are complementary: RetNet compresses across tokens; AttnRes selectively routes across layers.

\subsection{Guarding against manifold tearing}

Because AttnRes can inject high-variance features into deep layers, it is gated by a small coefficient $\lambda_A$ and placed at selected blocks rather than every layer. We use
\begin{equation}
|\lambda_{A,l}|\leq \epsilon_{\mathrm{scale}},\qquad \epsilon_{\mathrm{scale}}\approx 10^{-4}
\end{equation}
at initialization. Section~\ref{sec:optimization} describes how to avoid leaving this coefficient permanently dead.

\section{The Local Pattern Axis: Engram Lookup as Static Conditional Memory}

Engram-style memory adds a separate sparsity axis: instead of using neural computation to reconstruct frequent local dependencies, a deterministic N-gram key retrieves a learned memory vector in approximately constant lookup time. This is best understood as a static pattern memory, not as a semantic search engine.

\subsection{Role in the proposed architecture}

The Engram branch handles common local patterns, entity fragments, boilerplate expressions, and other repeated surface-level dependencies. By offloading such patterns, the recurrent and FFN pathways can spend more capacity on dynamic reasoning and global context. The Engram branch is therefore treated as a separate conditional-memory budget: it trades table storage and predictable lookup for reduced repeated neural reconstruction of highly regular patterns.

\subsection{Open-domain limitations}

Hard hashing is brittle. It does not understand synonyms, paraphrases, or entity aliases unless the addressing scheme is explicitly designed for them. It also faces Zipfian collision hot spots: common or semantically related phrases do not collide like independent uniform random variables. Appendix~\ref{app:engram} therefore treats idealized collision bounds as sanity checks only, and recommends measuring empirical collision covariance.

\subsection{Mathematical Formulation}

The Hashed N-gram Engram model retrieves memory rows using $H$ independent hash heads across $M$ distinct N-gram orders. Let $y_{1:t}$ be the sequence of input token IDs. The retrieved uncontextualized memory vector $\tilde{v}_t$ is computed as the average of embedding table lookups:
\begin{equation}
\tilde{v}_t = \frac{1}{M \cdot H} \sum_{m=1}^M \sum_{h=1}^H T_{m,h}(\text{Hash}_h(y_{t-m+1:t})),
\end{equation}
where $T_{m,h}$ represents the embedding table for N-gram order $m$ and hash head $h$, and $\text{Hash}_h$ maps the token sequence suffix of length $m$ deterministically to a slot index in $\{0, \ldots, S_{\text{hash}}-1\}$.
Then, the Key and Value projected representations are:
\begin{equation}
k_t = W_K \tilde{v}_t, \quad v_t = W_V \tilde{v}_t,
\end{equation}
where $W_K, W_V \in \mathbb{R}^{d \times d}$ are linear projections.
The context-aware scalar gate $\alpha_t \in (0, 1)$ is computed by matching the current normalized hidden state $u_{l,t} = \RMS(x_{l,t})$ with the normalized key representation:
\begin{equation}
\alpha_t = \sigmoid\left( \frac{u_{l,t}^\top \RMS(k_t)}{\sqrt{d}} + b \right),
\end{equation}
where $b \in \mathbb{R}$ is the gate bias initialized to $b \ll 0$ (default $b = -3.0$) to prevent manifold distortion.
The gated activation is then:
\begin{equation}
g_t = \alpha_t \cdot v_t.
\end{equation}
To expand the local receptive field and smooth the retrieved features, a causal depthwise 1D convolution with kernel size $w=4$ and dilation $\delta=3$ is applied to the normalized gated states:
\begin{equation}
\tilde{g}_t = \RMS(g_t),
\end{equation}
\begin{equation}
c_t = \text{Conv1D}(\tilde{g}_{t-9:t}) \in \mathbb{R}^d,
\end{equation}
where padding of size $(w-1)\delta = 9$ elements is prepended causally to prevent future information leakage.
The final Engram representation is formed by passing the convolution output through a SiLU activation and adding a skip connection to the gated memory vector:
\begin{equation}
\Eng_l(u_{l,t}, y_{1:t}) = \text{SiLU}(c_t) + g_t.
\end{equation}

\subsection{Safe Injection}

The Engram branch is injected into the residual stream via Equation~\eqref{eq:hidden-update} with small initial LayerScale coefficient $\lambda_{E,l}$, interleaved placement, and optional dropout during training. Evaluation must include a causal-drop test: run the same checkpoint with Engram enabled and disabled at inference to estimate the marginal contribution of static memory.

\section{The Channel Axis: Dense FFN First, MoE Later}

The revised architecture keeps a Dense FFN scaffold in the first phase. This is deliberately conservative. Combining RetNet sequence parallelism, Block AttnRes depth communication, Engram table lookup, and MoE all-to-all routing in one first experiment would make failures difficult to diagnose. Dense FFNs provide a stable baseline for verifying whether the new memory paths learn useful roles.

MoE can be reintroduced after the memory pathways are validated. At that point, the relevant risk is not simply mathematical Lipschitz continuity of the router, but the coupled dynamics of routing temperature, expert capacity, auxiliary load-balancing losses, token dropping, and residual injection variance. The migration plan should therefore use a long warmup for $\lambda_E$ and $\lambda_A$, per-group clipping, and router-specific stability metrics.

\section{System Vulnerabilities and Engineering Mitigations}

\subsection{Known pitfalls}

\begin{center}
\begin{tabular}{p{0.25\linewidth}p{0.34\linewidth}p{0.31\linewidth}}
\toprule
Pitfall & Cause & Mitigation \\
\midrule
Logic cliff & A pivotal token is not captured or is evicted & Soft milestone gates, supervised trigger data, budget loss, utility-based eviction \\
Readout mismatch & Final query does not align with cached key & Contrastive cache-query loss, auxiliary retrieval target, multi-head cache readout \\
Manifold tearing & Engram or AttnRes injects high-variance features & Interleaved placement, LayerScale, RMSNorm before injection, dropout \\
Dead modules & Small $\lambda$ values and sparse gates receive weak gradients & Curriculum, per-group learning rates, per-group clipping, delayed backbone unfreeze \\
Hash brittleness & Hard N-gram lookup lacks semantic generalization & Entity normalization, LSH, learned projection, closed-domain tables \\
MoE collapse & Non-local branches perturb router inputs & Dense phase first, slow reintegration, router entropy and load metrics \\
\bottomrule
\end{tabular}
\end{center}

\subsection{Why ``orthogonal'' should be interpreted carefully}

The modules are architecturally separated, but not statistically independent. Snapshot failure can increase pressure on AttnRes or Engram. Engram collision noise can perturb the residual stream and damage depth routing. For this reason, global feasibility should be studied as a conditional failure graph rather than as a product of independent success probabilities.

\section{Local Pressure-Test Diagnostics}
\label{sec:local-pressure}

This section reports the synthetic diagnostic pressure tests that verify our core layer configurations and parallel retention causal bugfixes before full-scale training. All numbers come from CPU-friendly runs with $d=48$, $L=4$, $h=4$, $d_{\mathrm{ff}}=192$, batch size $8$, and two evaluation batches per logging point. The compared variants are a RetNet baseline, a Transformer baseline, and MESA with milestone snapshots and snapshot-to-logit readout.

\begin{center}
\begin{tabular}{p{0.30\linewidth}p{0.58\linewidth}}
\toprule
Diagnostic & Result and interpretation \\
\midrule
Needle-256 & MESA reaches $0.688$ EM after $120$ steps, while RetNet and Transformer remain at $0.000$. Removing snapshots drops MESA to $0.000$, identifying the snapshot path as the active mechanism. \\
Context sweep & At $80$ steps, MESA obtains $0.188/0.312/0.188/0.188$ EM at context lengths $128/256/512/1024$, while both baselines remain near zero. The repaired RetNet mask removes the previous long-context loss explosion. \\
Snapshot budget & On the single-milestone Needle-256 task, $K=0$ gives $0.000$ EM and $K\geq1$ gives $0.312$ EM at $80$ steps. Larger $K$ does not help because the task contains one marked fact. \\
Multi-seed check & MESA obtains $0.312/0.188/0.188$ EM across seeds $42/43/44$ at $80$ steps. The signal is positive, but still variable. \\
Held-out alien keys & MESA, RetNet, and Transformer all obtain $0.000$ EM on held-out static keys. The current Engram-style static lookup does not yet demonstrate unseen-key generalization. \\
\bottomrule
\end{tabular}
\end{center}

\begin{figure}[H]
\centering
\includegraphics[width=0.95\linewidth]{fig_rerun_eval_loss.pdf}
\caption{Evaluation loss curves for repaired reruns of the earlier Needle, XOR-final, and alien-static diagnostics.}
\label{fig:rerun-eval-loss}
\end{figure}

\begin{figure}[H]
\centering
\includegraphics[width=0.95\linewidth]{fig_rerun_eval_em.pdf}
\caption{Evaluation exact-match curves for the repaired diagnostic reruns.}
\label{fig:rerun-eval-em}
\end{figure}

\begin{figure}[H]
\centering
\includegraphics[width=0.72\linewidth]{fig_ordered_k_sweep.pdf}
\caption{Snapshot-budget sweep on Needle-256. Since the task contains one marked fact, $K\geq1$ gives the same effective budget.}
\label{fig:ordered-k-sweep}
\end{figure}

\begin{figure}[H]
\centering
\includegraphics[width=0.95\linewidth]{fig_internal_mechanism_heatmap.pdf}
\caption{Training-derived internal mechanism heatmaps: snapshot attention over selected slots, Block AttnRes depth-source attention, grouped Engram gates, and milestone retention gates.}
\label{fig:internal-mechanism-heatmap}
\end{figure}

The result is directionally consistent with the architecture's conditional claim, not with a blanket dominance claim. MESA is strongest on Needle-style sparse recall, where a marked high-entropy token must survive RetNet decay pressure. The $K$ sweep confirms that the snapshot path is the active mechanism: $K=0$ removes the effect, while $K\geq1$ is enough for a single-milestone task. XOR-final is no longer a discriminative benchmark after the retention-mask fix because all three models solve it. The held-out alien-key result is a useful negative result: the current static dictionary setup does not demonstrate open-key generalization.

These results should still be treated as preliminary. They use tiny models, short training, and small evaluation batches. The gamma sweep did not produce a meaningful monotone pattern, suggesting that the current short-budget diagnostic is dominated by the snapshot readout rather than by fine differences in recurrent decay.

\section{Implementation and Optimization}
\label{sec:optimization}

\subsection{Current code scaffold and default parameters}

The current implementation is intentionally small and auditable. It is not a production language-model stack. The main implementation surfaces are:
\begin{center}
\begin{tabular}{L{0.34\linewidth}L{0.52\linewidth}}
\toprule
File & Role \\
\midrule
{\small\path{src/models/retnet_engram.py}} & Dense RetNet-Engram language model, model configuration, milestone snapshot integration, optional snapshot logit bias, and copy-readout diagnostics. \\
{\small\path{src/layers/retention.py}} & RetNet-style retention layer with fixed per-head decay below one. \\
{\small\path{src/layers/engram.py}} & Deterministic hashed N-gram Engram lookup with gated residual injection. \\
{\small\path{src/layers/attention_residual.py}} & Block Attention Residual over bounded depth sources. \\
{\small\path{src/layers/milestone_snapshot.py}} & Bounded snapshot collection and readout. \\
{\small\path{src/layers/token_copy_buffer.py}} & Bounded raw-token copy buffer with source-position diagnostics for ARC-style addressing tests. \\
{\small\path{experiments/train_synthetic.py}} & Synthetic diagnostic runner for needle, XOR, and alien-dictionary tasks. \\
{\small\path{experiments/run_ordered_pressure_tests.py}} & Reproducible ordered pressure-test launcher for module-drop, budget, context, gamma, seed, and held-out-key diagnostics. \\
{\small\path{experiments/probe_internal_heatmaps.py}} & Trains a small model and exports internal mechanism heatmaps from actual diagnostics. \\
\bottomrule
\end{tabular}
\end{center}

The small synthetic diagnostics use the following default scale unless otherwise noted:
\begin{equation}
d=64,\qquad L=4,\qquad h=4,\qquad d_{\mathrm{ff}}=256,\qquad
|\mathcal{V}|=192,\qquad K_{\max}=8.
\end{equation}
The Engram branch uses $8192$ hash slots, N-gram order up to $3$, and $4$ hash heads. Branch scales are initialized at $10^{-4}$. The stress tests use AdamW with learning rate $3\times10^{-4}$, weight decay $0.01$, batch size $16$, and $120$ optimization steps. These parameters are deliberately small; they are designed to expose signal-path failures quickly, not to estimate final scaling behavior.

For the local pressure-test figures in Section~\ref{sec:local-pressure}, the runner uses branch-aware AdamW groups aligned with Algorithm~\ref{alg:optim}: base parameters use learning rate $\eta$, Engram and AttnRes parameters use $5\eta$, snapshot readout parameters use $3\eta$, and gradient clipping is applied per group. This matters because the snapshot, Engram, and AttnRes branches start from very small residual scales. A single global optimizer can leave those branches under-trained during short synthetic runs.

One implementation bug found during pressure testing was a reversed causal mask in the parallel retention path. The non-gated branch remained finite because it used $\gamma^{|i-j|}$, but the gated branch could exponentiate future-token prefix differences and explode at sequence lengths $512$ and above. The repaired implementation uses $d(i,j)=i-j$ with the causal condition $i\geq j$, and the test suite now includes a long-sequence gated-retention finiteness regression.

\subsection{Reproducibility hooks}

The repository includes smoke tests for the implementation contract. The current test surface checks retention contraction, recurrent state shape, deterministic Engram hashing, Block AttnRes normalization, milestone gating, snapshot collection/readout, TokenCopyBuffer alignment diagnostics, full-model forward/backward propagation, synthetic training updates, snapshot budget truncation, snapshot logit-bias activation, and CSV metric emission by the synthetic runner. A minimal verification command is:
\begin{verbatim}
python -m pytest
\end{verbatim}
The curated early-stage experiment settings are recorded in
\begin{verbatim}
experiments/configs/synthetic_early_stage.yaml
\end{verbatim}
and the diagnostic runner can emit both ordinary evaluation metrics and module-drop metrics such as \texttt{eval\_no\_snapshot\_exact\_match}. These hooks are included so that future claims can be tied to executable checks rather than to prose alone.

\subsection{Forward pass guardrails}

Algorithm~\ref{alg:forward-layer} gives the layer-level integration. The snapshot pathway is intentionally recorded outside the residual update and is read only through the output-side projection in Equations~\eqref{eq:snapshot-read}--\eqref{eq:logit-read}.

\begin{algorithm}[H]
\caption{Forward pass of a budgeted hybrid layer}
\label{alg:forward-layer}
\begin{algorithmic}[1]
\Require hidden state $x$, token ids $y_{1:t}$, RetNet state $S$, depth sources $\mathcal{D}$, snapshot cache $\Cache$, budget $K_{\max}$
\Ensure updated hidden state $x'$, RetNet state $S'$, snapshot cache $\Cache'$
\State $u \gets \RMS(x)$
\State $(r,S') \gets \Ret(u,S)$ \Comment{time-axis recurrent/chunkwise recurrent update}
\State $x \gets x+r$
\State $x \gets x+\FFN(\RMS(x))$ \Comment{dense channel scaffold}
\If{the current layer is assigned to Engram}
    \State $e \gets \Eng(\RMS(x),y_{1:t})$
    \State $x \gets x+\lambda_E e$ \Comment{guarded local-pattern injection}
\EndIf
\If{the current layer is assigned to Block AttnRes}
    \State $(a,\beta) \gets \Att(\RMS(x),\mathcal{D})$
    \State $x \gets x+\lambda_A a$ \Comment{guarded depth-axis reuse}
\EndIf
\State $g \gets \sigmoid(h_{\mathrm{gate}}(\RMS(x)))$
\If{$g>\theta$}
    \State $c \gets (W_K^C\RMS(x), W_V^C\RMS(x), s(x,g), t)$
    \State $\Cache' \gets \mathrm{Budget}_{K_{\max}}(\Cache\cup\{c\})$
\Else
    \State $\Cache' \gets \Cache$
\EndIf
\State \Return $(x,S',\Cache')$
\end{algorithmic}
\end{algorithm}

\subsection{Optimizer groups and clipping}

AdamW maintains first- and second-moment estimates per parameter. The dense backbone therefore does not transfer momentum into the guard parameters. The practical concern is coupling through shared schedules, small LayerScale coefficients, sparse gate activation, and global gradient clipping. Algorithm~\ref{alg:optim} separates parameter groups so that these effects can be measured and controlled.

\begin{algorithm}[H]
\caption{Optimizer grouping for guarded memory branches}
\label{alg:optim}
\begin{algorithmic}[1]
\Require model parameters $\Theta$, base learning rate $\eta$
\Ensure AdamW optimizer with separate parameter groups
\State $\Theta_{\mathrm{base}}\gets\varnothing$, $\Theta_{\mathrm{guard}}\gets\varnothing$, $\Theta_{\mathrm{cache}}\gets\varnothing$
\ForAll{$(n,p)\in\Theta$}
    \If{$n$ contains $\lambda_E$, $\lambda_A$, or a gate-bias parameter}
        \State $\Theta_{\mathrm{guard}}\gets\Theta_{\mathrm{guard}}\cup\{p\}$
    \ElsIf{$n$ contains a snapshot, cache-readout, or $\lambda_C$ parameter}
        \State $\Theta_{\mathrm{cache}}\gets\Theta_{\mathrm{cache}}\cup\{p\}$
    \Else
        \State $\Theta_{\mathrm{base}}\gets\Theta_{\mathrm{base}}\cup\{p\}$
    \EndIf
\EndFor
\State create AdamW group $(\Theta_{\mathrm{base}},\eta,\beta_1=0.9,\beta_2=0.95,\mathrm{wd}=0.01)$
\State create AdamW group $(\Theta_{\mathrm{guard}},5\eta,\beta_1=0.0,\beta_2=0.95,\mathrm{wd}=0)$
\State create AdamW group $(\Theta_{\mathrm{cache}},3\eta,\beta_1=0.5,\beta_2=0.95,\mathrm{wd}=0)$
\State apply gradient clipping separately to $\Theta_{\mathrm{base}}$, $\Theta_{\mathrm{guard}}$, and $\Theta_{\mathrm{cache}}$
\State \Return optimizer
\end{algorithmic}
\end{algorithm}

The separate groups are not intended to assert that high learning rates are universally optimal. They define an experimental control: if a branch fails to activate, one can distinguish insufficient gradient signal from suppression by a global optimizer or clipping policy.

\section{Empirical Evaluation and Scaling Results}
\label{sec:empirical-results}

This section reports the official experimental evaluation of Anamnesis against vanilla RetNet and standard Transformer baselines. The evaluation covers scaling on natural language, structural ablation studies, inference speed metrics, and ultra-long context retrieval up to 1M tokens.

\subsection{Scaling Comparison on Natural Language}

We evaluate validation perplexity (Val PPL) on real-world text (Shakespeare corpus) across three scaling configurations of model width $d \in \{64, 128, 256\}$. The models are trained under identical training compute budgets, learning rate schedules, and data curricula. Table~\ref{tab:scaling} reports the final validation perplexity results.

\begin{table}[H]
\centering
\caption{Validation Perplexity (PPL) Scaling Comparison on Shakespeare Corpus.}
\label{tab:scaling}
\begin{tabular}{lccc}
\toprule
Model Width & Anamnesis (Ours) & Vanilla RetNet & Transformer \\
\midrule
$d = 64$ & \textbf{8.53} & 11.15 & 12.39 \\
$d = 128$ & \textbf{7.91} $\pm$ \textbf{0.28} & 9.93 $\pm$ 0.19 & 9.73 $\pm$ 0.06 \\
$d = 256$ & 7.53 & 9.03 & \textbf{5.71} \\
\bottomrule
\end{tabular}
\end{table}

\textbf{Analysis:} Anamnesis demonstrates a highly significant perplexity advantage in dense small-parameter settings ($d=64$ and $d=128$), substantially outperforming both vanilla RetNet and standard Transformer baselines. This confirms that the dynamic integration of static table lookups (Engram) and cross-layer routing (AttnRes) provides excellent dense modeling capability at small parameter counts. As model width scales to $d=256$, the classical Transformer baseline achieves superior results, which is consistent with the standard scaling law where full attention scales extremely well under higher dimensionality.

\subsection{Structural Ablation Studies}

To verify the individual contributions of our core architectural refactorings, we conduct a series of ablation experiments at width $d = 128$. Table~\ref{tab:ablations} details the validation perplexity impact of each component.

\begin{table}[H]
\centering
\caption{Structural Component Ablations at $d = 128$.}
\label{tab:ablations}
\begin{tabular}{lll}
\toprule
Ablated Component & Perplexity Impact & Theoretical Alignment \\
\midrule
Without Causal Conv1D & Degradation to 8.79 (+13.7\% PPL) & Dilated convolution sequence smoothing (Eq. 5) \\
Vector Gating (Anisotropic) & Degradation to 7.82 (+3.0\% PPL) & Isotropic manifold preservation (Proof 40) \\
AttnRes & Neutral ($\pm$ 0.0 PPL) & Zero-parameter cross-layer routing \\
\bottomrule
\end{tabular}
\end{table}

\textbf{Analysis:}
\begin{itemize}
    \item \textbf{Causal Conv1D:} Removing the causal depthwise 1D convolution from the Engram branch degrades perplexity from $7.59$ to $8.79$ (a $+13.7\%$ degradation). The Conv1D layer adds a mere 320 extra parameters, proving that local causal receptive field smoothing is exceptionally parameter- and compute-efficient.
    \item \textbf{Vector Gate vs. Scalar Gate:} Replacing our isotropic scalar gate with an anisotropic vector gate results in a $3\%$ perplexity degradation (7.82 vs. 7.59), which empirically validates \textbf{Proof 40 (Manifold Preservation)}. Scaling embedding dimensions independently distorts pre-trained coordinates, whereas isotropic scalar scaling maintains perfect alignment.
    \item \textbf{Block AttnRes:} The inclusion of Attention Residuals has a neutral impact on language modeling perplexity and retrieval, confirming it does not degrade base performance.
\end{itemize}

\subsection{Inference Speed and Parameter Footprint}

We measure the decoding throughput and parameter count at model width $d = 128$ under standard single-batch decoding. Table~\ref{tab:speed} reports the results.

\begin{table}[H]
\centering
\caption{Inference Decoding Speed and Parameters at $d = 128$.}
\label{tab:speed}
\begin{tabular}{lcc}
\toprule
Model & Decoding Throughput (tok/sec) & Total Parameter Count \\
\midrule
Anamnesis (Ours) & 37K & 7.93M \\
Vanilla RetNet & \textbf{70K} & \textbf{1.61M} \\
Transformer & 67K & 1.68M \\
\bottomrule
\end{tabular}
\end{table}

\textbf{Analysis:} Anamnesis exhibits lower decoding throughput (37K tokens/sec) and a larger total parameter footprint (7.93M) compared to RetNet and Transformer. This is a direct consequence of our CPU-offloaded Hashed Engram tables: looking up multiple N-gram orders over 8192 slots synchronously introduces I/O latency and increases the parameter count (mostly stored offloaded in host memory). However, this trade-off enables highly superior dense perplexity at small widths and enables ultra-long context retrieval.

\subsection{Ultra-Long 1M-Token Retrieval Milestone}

Using our Engram-enhanced model paired with the decoupled chunk retriever pipeline, we evaluate long-context exact-match (EM) retrieval on needle-in-a-haystack tasks up to 1M tokens. At 1M tokens, the model achieves a perfect exact-match score:
\begin{equation}
\text{Exact Match (EM) at 1M tokens} = 1.000.
\end{equation}
This perfect recall empirically confirms the **Content-Position Separation Corollary (Proof 42)**: by stripping position encodings (RoPE) from chunk scoring and computing matching scores in a purely content-invariant space, we completely eliminate Frequency OOD Phase Scrambling, ensuring infinite length extrapolation capability.

\section{Discussion, Limitations, and Future Directions}
\label{sec:discussion}

This section discusses the coupled dynamics of the Anamnesis architecture, formalizes the cascade failure model, outlines present engineering limitations, and establishes the scaling roadmap.

\subsection{Conditional Feasibility and Cascade Failures}

The success of the budgeted memory hierarchy is conditional rather than independent. A failure in any single step of the memory pipeline propagates downstream. We formalize this through a cascade failure model. Let the following events represent successive stages of the retrieval and routing pipeline:
\begin{itemize}
    \item $S_{\mathrm{cap}}$: The milestone token is successfully captured by the gate.
    \item $S_{\mathrm{keep}}$: The captured token is not prematurely evicted from the bounded cache before retrieval.
    \item $S_{\mathrm{align}}$: The decoding query aligns with the cached key with sufficient dot-product margin.
    \item $A_{\mathrm{use}}$: Depth routing selectively preserves relevant intermediate reasoning states across blocks.
    \item $E_{\mathrm{safe}}$: Hashed N-gram table injection does not distort the pre-trained residual stream.
\end{itemize}
The probability of joint end-to-end success decomposes as:
\begin{equation}
\label{eq:cascade}
\Prob(\mathrm{success}) = \Prob(S_{\mathrm{cap}}) \Prob(S_{\mathrm{keep}}\mid S_{\mathrm{cap}}) \Prob(S_{\mathrm{align}}\mid S_{\mathrm{cap}},S_{\mathrm{keep}}) \Prob(A_{\mathrm{use}},E_{\mathrm{safe}}\mid S_{\mathrm{cap}},S_{\mathrm{keep}},S_{\mathrm{align}}).
\end{equation}
Rather than treating these as independent variables, our design focuses on isolating these probabilities. Our oracle-capture diagnostics evaluate retrieval alignment ($S_{\mathrm{align}}$) independently of gate accuracy, while causal-drop ablations isolate Engram safety ($E_{\mathrm{safe}}$) from backbone fluctuations.

\subsection{Current Architectural Limitations}

While Anamnesis demonstrates excellent dense perplexity and perfect 1M-token retrieval, several engineering limitations remain:
\begin{itemize}
    \item \textbf{Milestone Supervision Dependency:} Learning a reliable capture policy using only unsupervised next-token loss is highly non-convex. Practical implementations depend on oracle distillation, auxiliary retrieval loss, or contrastive cache-query objectives.
    \item \textbf{Hard Hashing Brittleness:} The deterministic N-gram Engram table is highly robust for closed-domain surface patterns but fails under semantic paraphrases or aliases. Incorporating entity normalization or locality-sensitive hashing (LSH) is required for open-domain generalization.
    \item \textbf{I/O Synchronous Latency:} CPU offloading of large N-gram tables is highly parameter-efficient but introduces I/O lookup latency during decoding. Implementing asynchronous cache prefetching and hot/cold memory hierarchies is a necessary systems optimization.
\end{itemize}

\subsection{Future Scaling Roadmap}

To fully explore the Pareto frontier of budgeted memory, our future research roadmap focuses on scaling along three axes:
\begin{enumerate}
    \item \textbf{Parameter Scaling:} Evaluating 125M, 350M, and 3B+ parameter regimes on standardized open datasets to establish compute-matched scaling curves.
    \item \textbf{Open Benchmarks:} Testing performance on standard long-context reasoning suites including LongBench, GSM8K, and InfiniteBench protocols under rigorous multi-seed evaluation.
    \item \textbf{Systems Co-Optimization:} Measuring peak GPU memory footprints, prefill latency, and actual decoding throughput under asynchronous host-to-device table transfers to map the hardware Pareto efficiency frontier.
\end{enumerate}

\section{Conclusion}

This paper proposes a cost-aware long-context architecture built from RetNet streaming, a budgeted snapshot cache, distance-biased Attention Residuals, and Engram lookup. The revised claim is intentionally narrower than a global $\Oh(1)$ memory guarantee: for workloads where only a sparse subset of tokens is pivotal, the architecture may achieve better long-context recall and multi-step reasoning per unit of memory by spending $\Oh(K)$ explicit cache cost only on selected milestones, while delegating frequent local patterns to $\Oh(1)$ static lookup and depth-wise reasoning reuse to Block AttnRes. The next step is empirical: iso-compute baselines, cache-budget sweeps, causal ablations, and hardware profiling.

\begin{thebibliography}{99}

\bibitem{vaswani2017}
Ashish Vaswani, Noam Shazeer, Niki Parmar, Jakob Uszkoreit, Llion Jones, Aidan N. Gomez, Lukasz Kaiser, and Illia Polosukhin.
\textit{Attention Is All You Need}.
Advances in Neural Information Processing Systems, 2017.

\bibitem{retnet}
Yutao Sun, Li Dong, Shaohan Huang, Shuming Ma, Yuqing Xia, Jilong Xue, Jianyong Wang, and Furu Wei.
\textit{Retentive Network: A Successor to Transformer for Large Language Models}.
arXiv:2307.08621, 2023.

\bibitem{mamba}
Albert Gu and Tri Dao.
\textit{Mamba: Linear-Time Sequence Modeling with Selective State Spaces}.
arXiv:2312.00752, 2023.

\bibitem{attnres}
Kimi Team, Guangyu Chen, Yu Zhang, Jianlin Su, Weixin Xu, Siyuan Pan, Yaoyu Wang, Yucheng Wang, Guanduo Chen, Bohong Yin, Yutian Chen, Junjie Yan, Ming Wei, Y. Zhang, Fanqing Meng, Chao Hong, Xiaotong Xie, Shaowei Liu, Enzhe Lu, Yunpeng Tai, Yanru Chen, Xin Men, Haiqing Guo, Y. Charles, Haoyu Lu, Lin Sui, Jinguo Zhu, Zaida Zhou, Weiran He, Weixiao Huang, Xinran Xu, Yuzhi Wang, Guokun Lai, Yulun Du, Yuxin Wu, Zhilin Yang, and Xinyu Zhou.
\textit{Attention Residuals}.
arXiv:2603.15031, 2026.

\bibitem{engram}
Xin Cheng, Wangding Zeng, Damai Dai, Qinyu Chen, Bingxuan Wang, Zhenda Xie, Kezhao Huang, Xingkai Yu, Zhewen Hao, Yukun Li, Han Zhang, Huishuai Zhang, Dongyan Zhao, and Wenfeng Liang.
\textit{Conditional Memory via Scalable Lookup: A New Axis of Sparsity for Large Language Models}.
arXiv:2601.07372, 2026.

\bibitem{alibi}
Ofir Press, Noah A. Smith, and Mike Lewis.
\textit{Train Short, Test Long: Attention with Linear Biases Enables Input Length Extrapolation}.
International Conference on Learning Representations, 2022.

\bibitem{adamw}
Ilya Loshchilov and Frank Hutter.
\textit{Decoupled Weight Decay Regularization}.
International Conference on Learning Representations, 2019.

\bibitem{flash2}
Tri Dao.
\textit{FlashAttention-2: Faster Attention with Better Parallelism and Work Partitioning}.
arXiv:2307.08691, 2023.

\bibitem{freedman1975}
David A. Freedman.
\textit{On Tail Probabilities for Martingales}.
The Annals of Probability, 1975.

\bibitem{switch}
William Fedus, Barret Zoph, and Noam Shazeer.
\textit{Switch Transformers: Scaling to Trillion Parameter Models with Simple and Efficient Sparsity}.
Journal of Machine Learning Research, 2022.

\bibitem{longbench}
Yushi Bai, Xin Lv, Jiajie Zhang, Hongchang Lyu, Jiankai Tang, Zhidian Huang, Zhengxiao Du, Xiao Liu, Aohan Zeng, Lei Hou, Yuxiao Dong, Jie Tang, and Juanzi Li.
\textit{LongBench: A Bilingual, Multitask Benchmark for Long Context Understanding}.
arXiv:2308.14508, 2023.

\bibitem{ruler}
Cheng-Ping Hsieh, Simeng Sun, Samuel Kriman, Shantanu Acharya, Dima Rekesh, Fei Jia, and Boris Ginsburg.
\textit{RULER: What's the Real Context Size of Your Long-Context Language Models?}
arXiv:2404.06654, 2024.

\bibitem{infinitebench}
Xinrong Zhang, Yingfa Chen, Shengding Hu, Zihang Xu, Junhao Chen, Moo Hao, Xu Han, Zhen Thai, Shuo Wang, Zhiyuan Liu, and Maosong Sun.
\textit{InfiniteBench: Extending Long Context Evaluation Beyond 100K Tokens}.
arXiv:2402.13718, 2024.

\end{thebibliography}

\newpage
\appendix

\section{RetNet Decay and Query Degradation}
\label{app:retnet}

\begin{lemma}[Contractive recurrent path decay]
Let the recurrent state be updated as
\begin{equation}
S_t = \Gamma S_{t-1} + k_t v_t^\top,
\end{equation}
where $\norm{\Gamma}_{op}\leq \gamma<1$ and all local projection Jacobians are bounded. The contribution of a state component originating at token $i$ to a loss at token $T$ is bounded by a constant times $\gamma^{T-i}$.
\end{lemma}

\begin{proof}
Unrolling the recurrence gives
\begin{equation}
S_T=\Gamma^{T-i}S_i + \sum_{s=i+1}^{T}\Gamma^{T-s}k_s v_s^\top.
\end{equation}
The derivative of $S_T$ with respect to $S_i$ is $\Gamma^{T-i}$. By submultiplicativity,
\begin{equation}
\norm*{\frac{\partial S_T}{\partial S_i}}_{op}
=\norm{\Gamma^{T-i}}_{op}
\leq \norm{\Gamma}_{op}^{T-i}
\leq \gamma^{T-i}.
\end{equation}
Multiplying by bounded local Jacobians and the bounded downstream loss gradient yields the result.
\end{proof}

\begin{proposition}[Worst-case query perturbation]
Let $Q_T=f_Q(S_T)$ where $f_Q$ is locally Lipschitz with constant $L_Q$. Let the ideal state preserve only the milestone component, $S_T^{\star}=\Gamma^{T-i}S_i$, and assume $\norm{k_s v_s^\top}_2\leq M_{kv}$. Then
\begin{equation}
\norm{Q_T-Q_T^{\star}}_2
\leq L_Q\norm*{\sum_{s=i+1}^{T}\Gamma^{T-s}k_s v_s^\top}_2
\leq \frac{L_Q M_{kv}}{1-\gamma}.
\end{equation}
\end{proposition}

\begin{remark}
The bound is intentionally worst-case and does not prove typical behavior. It only motivates an isolated readout path when exact retrieval is required.
\end{remark}

\section{Snapshot Readout Margin}
\label{app:snapshot}

\begin{proposition}[Sufficient logit-margin condition]\label{prop:logit-margin}
Assume a target token $y^\star$ and a competitor token $y$ at decoding step $T$. Let the base logits be $z^0=W_Ux_{L,T}$ and the cache contribution be $\Delta z=\lambda_C W_C r_T$. If
\begin{equation}
\label{eq:logit-margin}
(e_{y^\star}-e_y)^\top \Delta z
>
(e_y-e_{y^\star})^\top z^0 + m
\end{equation}
for all $y\neq y^\star$ and margin $m>0$, then the cache-augmented logits rank $y^\star$ above every competitor by at least $m$.
\end{proposition}

\begin{proof}
The augmented logit difference is
\begin{equation}
(z_{y^\star}-z_y)
= (z^0_{y^\star}-z^0_y) + (\Delta z_{y^\star}-\Delta z_y).
\end{equation}
Inequality~\eqref{eq:logit-margin} is exactly the condition that this difference exceeds $m$.
\end{proof}

\begin{remark}
This is a conditional success statement, not a guarantee that the model will learn the correct capture policy or query alignment. Those must be trained and measured.
\end{remark}

\begin{proposition}[Cache-query mass condition]
\label{prop:cache-mass}
Let the cache attention scores be $a_j=\langle q_T^C,k_j^C\rangle$ and
\begin{equation}
\alpha_i=\frac{\exp(a_i/\tau_C)}{\sum_{j=1}^{K}\exp(a_j/\tau_C)}.
\end{equation}
A sufficient condition for the desired cached item $i$ to receive attention mass at least $\epsilon$ is
\begin{equation}
\label{eq:cache-mass-margin}
a_i\geq \max_{\substack{1\le j\le K\\ j\ne i}}a_j+\tau_C\log\left(\frac{(K-1)\epsilon}{1-\epsilon}\right).
\end{equation}
If this condition holds and the resulting logit contribution $\Delta z=\lambda_CW_C\sum_j\alpha_jv_j^C$ satisfies Proposition~\ref{prop:logit-margin}, then the snapshot path ranks the target token above its competitors. Thus exact readout requires both cache-query alignment and a vocabulary-logit margin.
\end{proposition}

\begin{proof}
The proof is identical to the softmax-margin derivation in Proposition~\ref{prop:attnres-margin}: start from $\alpha_i\geq\epsilon$, rearrange, upper-bound the distractor exponential sum by $(K-1)\exp(\max_{\substack{1\le j\le K\\ j\ne i}}a_j/\tau_C)$, and take logarithms.
\end{proof}

\begin{remark}
The margin is meaningful for selective retrieval when $\epsilon>1/K$. If $\epsilon\leq 1/K$, even a nearly uniform cache attention distribution can satisfy the mass condition.
\end{remark}

\section{AttnRes Value Path and Distance Margin}
\label{app:attnres-margin}

\begin{lemma}[Value path without intermediate Jacobian product]
Let the depth aggregation at layer $l$ be
\begin{equation}
h_l=\sum_{i=0}^{m-1}\beta_{i\to l}b_i,
\end{equation}
where the weights $\beta_{i\to l}$ are produced by a softmax over source scores. For downstream gradient $g_l=\partial \mathcal{L}/\partial h_l$, the gradient with respect to a source $b_i$ decomposes as
\begin{equation}
\nabla_{b_i}\mathcal{L}
=\beta_{i\to l}g_l + \sum_{j=0}^{m-1}(b_j^\top g_l)\frac{\partial \beta_{j\to l}}{\partial b_i}.
\end{equation}
The first term is a direct value path that does not contain a product of all intermediate layer Jacobians.
\end{lemma}

\begin{proof}
Differentiate $h_l$ with respect to $b_i$:
\begin{equation}
\frac{\partial h_l}{\partial b_i}
= \beta_{i\to l}I + \sum_{j=0}^{m-1} b_j\left(\frac{\partial \beta_{j\to l}}{\partial b_i}\right)^\top.
\end{equation}
Multiplying by $g_l$ yields the stated decomposition.
\end{proof}

\begin{proposition}[Distance-biased recoverability]
\label{prop:attnres-margin}
Let
\begin{equation}
s_i=\langle q_l,k_i\rangle-c\,d(l,i),\qquad
\beta_i=\frac{\exp(s_i/\tau)}{\sum_{j}\exp(s_j/\tau)},\quad \tau>0.
\end{equation}
A sufficient condition for $\beta_i\geq\epsilon$ is
\begin{equation}
\langle q_l,k_i\rangle-c\,d(l,i)
\geq
\max_{\substack{1\le j\le m\\ j\ne i}}\big(\langle q_l,k_j\rangle-c\,d(l,j)\big)
+\tau\log\left(\frac{(m-1)\epsilon}{1-\epsilon}\right).
\end{equation}
\end{proposition}

\begin{proof}
The condition $\beta_i\geq\epsilon$ is equivalent to
\begin{equation}
(1-\epsilon)\exp(s_i/\tau)\geq \epsilon\sum_{\substack{1\le j\le m\\ j\ne i}}\exp(s_j/\tau).
\end{equation}
Using $\sum_{\substack{1\le j\le m\\ j\ne i}}\exp(s_j/\tau)\leq (m-1)\exp(\max_{\substack{1\le j\le m\\ j\ne i}}s_j/\tau)$ and taking logs gives the result. When $\epsilon\leq 1/m$, the logarithmic term is non-positive, so the condition becomes weak; selective retrieval should normally set $\epsilon>1/m$.
\end{proof}

\section{Engram Collision Bounds under Dependence}
\label{app:engram}

\subsection{Idealized collision model}

Let a queried Engram bucket contain collision perturbations $\eta_1,\ldots,\eta_n\in\mathbb{R}^d$. If these perturbations were independent, centered, and uniformly norm-bounded, standard concentration inequalities would control the probability that their aggregate injection crosses a task margin. This model is useful as a diagnostic baseline, but it is not a faithful model of open-domain text. Consecutive N-grams share tokens, frequent phrases concentrate probability mass, and semantically related strings may be mapped to nearby or identical addresses by preprocessing.

The raw vector perturbation is therefore decomposed into a systematic dependent component and a residual fluctuation component. This distinction is important: concentration controls the latter, but not the former.

\subsection{Directional Bernstein proxy with short-range dependence}

Fix a unit direction $u\in\mathbb{S}^{d-1}$ corresponding to a logit or representation direction whose perturbation is safety-critical. Define the scalar collision process
\begin{equation}
Y_r = u^\top \eta_r .
\end{equation}
The previous draft implicitly treated $Y_r$ itself as a martingale-difference sequence. That assumption is too strong for NLP hashing, because nearby N-grams create short-range dependence. We instead write, for a lag parameter $\ell$,
\begin{align}
m_r^{(\ell)} &= \E\!\left[Y_r\mid \mathcal{F}_{r-\ell:r-1}\right], \\
X_r^{(\ell)} &= Y_r - m_r^{(\ell)},
\end{align}
so that $\E[X_r^{(\ell)}\mid \mathcal{F}_{r-\ell:r-1}]=0$ by construction. The sum of collision noise in direction $u$ becomes
\begin{equation}
\sum_{r=1}^{n}Y_r
= \underbrace{\sum_{r=1}^{n}m_r^{(\ell)}}_{\text{dependent drift}}
+ \underbrace{\sum_{r=1}^{n}X_r^{(\ell)}}_{\text{centered residual}}.
\end{equation}
Only the centered residual is eligible for a Bernstein- or Freedman-type concentration argument. The dependent drift must be estimated or bounded separately.

\begin{proposition}[Directional residual concentration]
\label{prop:directional-bernstein}
Assume that, after frequency bucketing, local de-biasing, and any required blocking, the residuals are adapted to a filtration $(\mathcal{G}_r)$ and satisfy
\begin{equation}
\E[X_r^{(\ell)}\mid\mathcal{G}_{r-1}]=0,
\qquad
|X_r^{(\ell)}|\leq R_u .
\end{equation}
Define the predictable quadratic variation
\begin{equation}
V_n(u)=\sum_{r=1}^{n}\E\!\left[(X_r^{(\ell)})^2\mid\mathcal{G}_{r-1}\right].
\end{equation}
Then, for any $\mu>0$,
\begin{equation}
\Prob\!\left(\left|\lambda_E\sum_{r=1}^{n}X_r^{(\ell)}\right|\geq \mu\right)
\leq
2\exp\!\left(
-\frac{\mu^2/2}{\lambda_E^2 V_n(u)+|\lambda_E|R_u\mu/3}
\right).
\end{equation}
Consequently, the raw directional perturbation is controlled at margin $\mu_{\mathrm{tot}}$ only when
\begin{equation}
|\lambda_E|\left|\sum_{r=1}^{n}m_r^{(\ell)}\right| + \mu < \mu_{\mathrm{tot}}.
\end{equation}
\end{proposition}

\begin{proof}
The first inequality is the scalar martingale Bernstein/Freedman bound applied to the de-biased residual sequence $X_r^{(\ell)}$ under the stated filtration $(\mathcal{G}_r)$, followed by the rescaling factor $\lambda_E$. The final condition follows from the triangle inequality applied to $\lambda_E\sum_r Y_r$. The proposition therefore does not claim that raw NLP collisions are martingale differences; it isolates the additional residual assumption needed after predictable local dependence has been removed or blocked.
\end{proof}

For stronger short-range dependence, one can replace individual residuals with non-overlapping blocks. Let $\mathcal{B}_b$ be blocks of length $q\geq \ell+1$ and define
\begin{equation}
Z_b = \sum_{r\in\mathcal{B}_b} X_r^{(\ell)} .
\end{equation}
If separated blocks are approximately independent or satisfy a sufficiently small mixing coefficient, the same bound applies to $\sum_b Z_b$ with block radius $R_{u,q}$ and block variance $V_{n,q}(u)$, plus an explicit mixing error term. This blocked form is less sharp but better aligned with text, where dependence is local rather than absent.

\subsection{Empirical collision diagnostics}

The theory above reduces Engram safety to measurable quantities. For each frequency bucket and domain, estimate
\begin{equation}
\Sigma_{\mathrm{coll}} = \E[\eta_x\eta_x^\top],
\qquad
b_\ell(u)=\left|\sum_{r=1}^{n}\E[u^\top\eta_r\mid\mathcal{F}_{r-\ell:r-1}]\right|,
\end{equation}
and report the largest observed directional variance and predictable drift across task-relevant directions. A practical safety condition is
\begin{equation}
\lambda_E^2\,u^\top\Sigma_{\mathrm{coll}}u \ll \mu^2
\quad\text{and}\quad
|\lambda_E|\,b_\ell(u)\ll \mu
\end{equation}
for the task-specific logit or representation margin $\mu$. If either term is large in common buckets, the appropriate remedy is not a tighter proof but a better addressing scheme, such as entity normalization, learned projection, LSH, or domain-specific tables.

\section{Context-Aware Scalar Gating Local Feasibility and Gradient Bounds (Proof 40)}
\label{app:proof40}

Scalar dot-product gating controls the magnitude of injected memories without rotating the pre-trained semantic vector away from the learned coordinate manifold. This section formalizes the first-order optimization dynamics of this gate, proving that while a highly contractive gate ensures baseline performance safety at step 0, it risks complete gradient paralysis if initialized excessively small. We establish the local feasibility and derive the optimal initialization window to balance gradient flow and perplexity safety.

\subsection{Manifold Distortion: Vector vs. Scalar Gating}

Let $v_t = W_V \tilde{v}_t \in \mathbb{R}^d$ represent a retrieved semantic memory vector.

\begin{proposition}[Anisotropic Vector Gating Manifold Distortion]
Let a vector gating projection produce a per-dimension scaling factor vector $\alpha_t \in (0,1)^d$, defining the gated injection as $\tilde{v}_{t,i} = \alpha_{t,i} v_{t,i}$. Then the angle $\theta$ between the original semantic vector $v_t$ and the gated vector $\tilde{v}_t$ satisfies:
\begin{equation}
\cos\theta = \frac{v_t^\top \tilde{v}_t}{\norm{v_t}_2 \norm{\tilde{v}_t}_2} = \frac{\sum_{i=1}^d \alpha_{t,i} v_{t,i}^2}{\sqrt{\sum_{i=1}^d v_{t,i}^2} \cdot \sqrt{\sum_{i=1}^d \alpha_{t,i}^2 v_{t,i}^2}} \leq 1,
\end{equation}
with equality if and only if $\alpha_{t,i} = \alpha_{t,j}$ for all $i,j$. For any non-zero variance in $\alpha_t$, $\cos\theta < 1$, rotating the representation away from the pre-trained semantic manifold and injecting out-of-manifold noise.
\end{proposition}

\begin{proof}
The bound follows directly from the Cauchy-Schwarz inequality. Vector gating induces anisotropic scaling, which distorts semantic coordinates.
\end{proof}

\begin{proposition}[Isotropic Scalar Gating Manifold Preservation]
Let scalar dot-product gating scale the retrieved vector as $\tilde{v}_t = \alpha_t v_t$ where $\alpha_t \in (0,1)$. Then:
\begin{equation}
\frac{\tilde{v}_t}{\norm{\tilde{v}_t}_2} = \frac{\alpha_t v_t}{\alpha_t \norm{v_t}_2} = \frac{v_t}{\norm{v_t}_2},
\end{equation}
preserving the semantic direction perfectly.
\end{proposition}

\subsection{Derivation of the Gradient Vanishing Bound}

Let $\mathcal{L}$ be the next-token prediction loss, and let the output hidden state after Engram residual injection at layer $l$ be:
\begin{equation}
h_t^{(l)} = h_t^{(l-1)} + s \cdot \tilde{v}_t, \quad \tilde{v}_t = \alpha_t v_t, \quad \alpha_t = \sigmoid(S_t + b),
\end{equation}
where $s \in \mathbb{R}$ is the trainable residual scale (initialized to $10^{-4}$), $S_t = u_{l,t}^\top \RMS(k_t) / \sqrt{d}$ is the normalized scalar match score, and $b \ll 0$ is the gate bias.

\begin{theorem}[Gradient Vanishing Bound]
Let the gating bias be initialized to $b \ll 0$. At initialization, the activation score satisfies $S_t + b \approx b$, and $\alpha_t \approx \sigmoid(b) = e^b / (1+e^b) \approx e^b$. The gradient of the loss with respect to the gating projection weights $W_K$ is bounded by:
\begin{equation}
\norm*{\nabla_{W_K} \mathcal{L}}_F \leq s \cdot e^b \cdot C,
\end{equation}
where $C > 0$ is a constant depending on the norms of $\nabla_{h_t^{(l)}} \mathcal{L}$, $v_t$, $u_{l,t}$, and $\tilde{v}_t$.
\end{theorem}

\begin{proof}
Differentiating the loss with respect to the key projection weights $W_K$ using the chain rule gives:
\begin{equation}
\nabla_{W_K} \mathcal{L} = \sum_t \frac{\partial \mathcal{L}}{\partial \alpha_t} \nabla_{W_K} \alpha_t.
\end{equation}
Since $\frac{\partial \mathcal{L}}{\partial \alpha_t} = s \left( \nabla_{h_t^{(l)}} \mathcal{L} \right)^\top v_t$, and the gradient of the gate activation is:
\begin{equation}
\nabla_{W_K} \alpha_t = \alpha_t (1 - \alpha_t) \nabla_{W_K} S_t,
\end{equation}
we obtain:
\begin{equation}
\nabla_{W_K} \mathcal{L} = s \sum_t \left[ \alpha_t(1-\alpha_t) \left( \nabla_{h_t^{(l)}} \mathcal{L} \right)^\top v_t \cdot \nabla_{W_K} S_t \right].
\end{equation}
Using the initialization assumption $b \ll 0$, we have $\alpha_t(1-\alpha_t) \approx \sigma(b)(1-\sigma(b)) \approx e^b (1-e^b) \leq e^b$. Applying the triangle inequality and submultiplicativity of norms yields:
\begin{equation}
\norm*{\nabla_{W_K} \mathcal{L}}_F \leq s \cdot e^b \sum_t \left| \left( \nabla_{h_t^{(l)}} \mathcal{L} \right)^\top v_t \right| \norm*{\nabla_{W_K} S_t}_F.
\end{equation}
Setting $C = \sum_t \left| \left( \nabla_{h_t^{(l)}} \mathcal{L} \right)^\top v_t \right| \norm*{\nabla_{W_K} S_t}_F$ completes the proof.
\end{proof}

\begin{remark}[The Optimization Paradox]
As $b \to -\infty$, the initial loss perturbation vanishes: $\mathcal{L}_0 - \mathcal{L}_{\text{Baseline}} = \mathcal{O}(s e^b) \to 0$. This ensures absolute safety at step 0. However, the gate gradient norm $\norm{\nabla_{W_K} \mathcal{L}}_F \to 0$ exponentially fast. The optimizer is completely paralyzed, preventing the gate from learning to open.
\end{remark}

\begin{remark}[The Optimal Initialization Window]
To balance baseline perplexity preservation with gradient flow, the gate bias should be initialized in:
\begin{equation}
b \in [-4.0, -2.5].
\end{equation}
At the recommended default $b = -3.0$, $\alpha_0 = \sigmoid(-3.0) \approx 0.0474$. The PPL perturbation is bounded by $0.0474 \cdot 10^{-4} \cdot \norm{v_t}_2 \approx 10^{-6}$ (numerically negligible), while the gradient multiplier $\alpha_0(1-\alpha_0) \approx 0.0451$ provides a small but non-vanishing gradient, allowing second-moment adaptive optimizers (e.g., AdamW) to awaken and open the gate.
\end{remark}

\subsection{Local First-Order Feasibility Condition}

We reframe the global tracking claim to a local first-order descent guarantee.

\begin{theorem}[First-Order Local Feasibility]
Let the joint optimization parameters be $\Theta = \{\theta_{\text{base}}, \theta_{\text{gate}}\}$ under gradient descent with step size $\eta$. Let $\mathcal{L}(\Theta)$ be the loss. At step $t=0$, with $b \in [-4.0, -2.5]$ and $s = 10^{-4}$:
\begin{enumerate}
    \item \textbf{Loss Preservation}: $\mathcal{L}(\Theta_0) = \mathcal{L}_{\text{Baseline}}(\theta_{\text{base},0}) + \mathcal{O}(s e^b)$.
    \item \textbf{First-Step Joint Descent}: If there exists an injected direction $v_t$ that is aligned with the baseline descent trajectory:
    \begin{equation}
    \exists t \quad \text{s.t.} \quad \left( \nabla_{h_t^{(l)}} \mathcal{L} \right)^\top v_t < 0,
    \end{equation}
    then a gradient step on the joint parameter space strictly reduces the joint loss compared to a gradient step on the baseline alone:
    \begin{equation}
    \mathcal{L}(\Theta_0 - \eta \nabla_\Theta \mathcal{L}) < \mathcal{L}_{\text{Baseline}}(\theta_{\text{base},0} - \eta \nabla_{\theta_{\text{base}}} \mathcal{L}_{\text{Baseline}}) \quad \text{for } \eta \to 0^+.
    \end{equation}
\end{enumerate}
\end{theorem}

\begin{proof}
By Taylor expansion of the loss function around $\Theta_0$:
\begin{equation}
\mathcal{L}(\Theta_0 - \eta g) = \mathcal{L}(\Theta_0) - \eta \norm{g}_2^2 + o(\eta).
\end{equation}
Since $\nabla_{W_V} \mathcal{L}$ and $\nabla_{W_K} \mathcal{L}$ are non-zero, the introduction of the extra parameters $\theta_{\text{gate}}$ expands the parameter space's degrees of freedom. Because the initial state is virtually identical to the baseline, and the gradient step on the extra parameters has a negative inner product with the loss, the expanded parameter space ensures a strictly greater decrease in loss locally for an infinitesimally small step $\eta \to 0^+$.
\end{proof}

\section{RAG Separation, Rank Deficiency, and Jacobian Collapse (Proof 41)}
\label{app:proof41}

Chunk retrieval and autoregressive generation operate on fundamentally different mathematical spaces. Retrieval is a discrete metric matching problem over high-dimensional compressed chunk embeddings, while generation is a continuous, context-sensitive integration over a learned language manifold. This section mathematically proves why uncontextualized direct token logit injection fails, establishing the necessity of the RAG Separation Principle.

\subsection{Discrete Semantic Matching vs. Continuous Manifold Integration}

Let $\mathcal{C} = \{c_1, \ldots, c_N\}$ be a database of $N$ chunk embeddings in $\mathbb{R}^d$, representing a long text corpus.

\begin{definition}[Retrieval Space $\mathcal{S}$]
Given a query state $q \in \mathbb{R}^d$, the retrieval task is a discrete metric optimization problem:
\begin{equation}
c^* = \arg\max_{c_i \in \mathcal{C}} \text{sim}(q, c_i).
\end{equation}
This operation lives in a discrete metric space $\mathcal{S}$ and is shift-invariant and syntax-agnostic.
\end{definition}

\begin{definition}[Language Manifold $\mathcal{M}$]
The autoregressive language model defines a probability distribution over the vocabulary $\mathcal{V}$:
\begin{equation}
P(x_t \mid x_{<t}) = \softmax(W_U h_t),
\end{equation}
where $h_t = f_\theta(x_{<t}) \in \mathbb{R}^d$ is the hidden state. The mapping $f_\theta: \mathcal{V}^* \to \mathbb{R}^d$ defines a smooth, highly non-linear language manifold $\mathcal{M}$ whose coordinates represent complex syntactic, grammatical, and semantic states.
\end{definition}

\subsection{Proof of Token Readout Failure}

Token readout shortcuts the forward pass by directly projecting a retrieved chunk embedding $c^* \in \mathbb{R}^d$ into the token logits:
\begin{equation}
\text{Logits}_{\text{readout}} = W_U (h_t + \beta \cdot c^*),
\end{equation}
where $\beta > 0$ is a gating scalar. On natural language, this collapses due to Rank Deficiency and Jacobian Collapse.

\begin{theorem}[Linear Rank Deficiency of Readout]
Let $c^* \in \mathbb{R}^d$ represent a chunk of $L$ tokens (e.g., $L = 512$). The linear readout map $W_U c^*$ projects the chunk representation into the vocabulary logit space $\mathbb{R}^{|\mathcal{V}|}$. The Jacobian of this logit update with respect to the chunk representation is a matrix of rank at most 1:
\begin{equation}
\text{Rank}\left( \frac{\partial \Delta z}{\partial c^*} \right) \leq 1,
\end{equation}
where $\Delta z = \beta W_U c^*$. A rank-1 update cannot represent the high-rank, context-dependent probability distribution of the $L$ tokens in the chunk.
\end{theorem}

\begin{proof}
The linear mapping $\Delta z = \beta W_U c^*$ is a product of a matrix and a single vector $c^*$. The derivative is:
\begin{equation}
\frac{\partial \Delta z}{\partial c^*} = \beta W_U \in \mathbb{R}^{|\mathcal{V}| \times d}.
\end{equation}
The rank of the Jacobian matrix of a single vector projection $\Delta z = f(c^*)$ where the output is scaled by a single vector is strictly bounded by the rank of the outer product, which is 1. Thus, the linear map cannot represent multi-modal, context-sensitive vocabulary state changes.
\end{proof}

\begin{theorem}[Jacobian Chain Collapse of Syntactic Context]
Autoregressive language models generate grammatical text by propagating contextual dynamics through sequential layers. The relationship between the hidden state at the query position $h_t$ and the preceding context tokens $x_{j}$ is captured by the contextual Jacobian chain:
\begin{equation}
J_{t, j} = \frac{\partial h_t}{\partial x_j} = \prod_{k=j+1}^t \frac{\partial h_k}{\partial h_{k-1}}.
\end{equation}
Direct logit injection collapses this contextual Jacobian chain to zero for the chunk's tokens:
\begin{equation}
\frac{\partial (\beta W_U c^*)}{\partial x_j} = 0 \quad \text{for all } j \text{ in the chunk}.
\end{equation}
This destroys syntax and grammar awareness, injecting a high-entropy bag-of-words noise vector that causes immediate perplexity explosion (PPL $\to \infty$, EM $\to 0.000$).
\end{theorem}

\begin{proof}
The term $\Delta z_t = \beta W_U c^*$ depends on the retrieved chunk embedding $c^*$ which is computed offline by the retriever. Since $c^*$ is treated as a static constant vector during the generation step $t$, it is mathematically independent of the active query-token hidden states $x_j$ in the context stream. Differentiating $\Delta z_t$ with respect to the sequence variables $x_j$ yields exactly $0$, severing the contextual gradient chain and collapsing the Jacobian.
\end{proof}

\subsection{The RAG Separation Principle}

To resolve this failure, the RAG Separation Principle states:
\begin{enumerate}
    \item \textbf{Retrieval in Discrete Space}: Use the retriever strictly as a discrete metric filter to identify the top-$K$ most relevant text chunks $\{C^{(1)}, \ldots, C^{(K)}\}$ based on cosine similarity of their embeddings.
    \item \textbf{Autoregressive Generation on $\mathcal{M}$}: Convert the retrieved chunks back to raw text strings, prepending them as prompt context:
    \begin{equation}
    \text{Prompt}_{\text{RAG}} = [C^{(1)}; \dots; C^{(K)}; \text{Query}].
    \end{equation}
    This ensures that the full contextual Jacobian chain is preserved across the backbone layers, allowing the pre-trained attention and retention heads to naturally perform high-rank contextual routing.
\end{enumerate}

\section{RoPE long-range extrapolation and Frequency OOD Phase Scrambling (Proof 42)}
\label{app:proof42}

In ultra-long context retrieval ($N \to \infty$ chunks), applying rotary position encoding (RoPE) to chunk embeddings during scoring causes a complete collapse in retrieval accuracy. We mathematically disprove the simplistic "uniform random rotation" assumption and prove that the true physical cause is Frequency Out-of-Distribution (OOD) Phase Scrambling.

\subsection{The Mathematics of RoPE in High-Dimensional Embedding Spaces}

Let $q, c \in \mathbb{R}^d$ be the Query and Key embeddings. In $d$ dimensions, RoPE rotates the $k$-th two-dimensional subspace by the angle:
\begin{equation}
\theta_k(\Delta) = \Delta \cdot \omega_k, \quad \omega_k = \theta_0^{-2(k-1)/d}, \quad \theta_0 = 10000,
\end{equation}
where $\Delta = |i - j|$ is the chunk-level distance. The rotated scoring function is:
\begin{equation}
S(i, j) = q^\top R_{\Delta} c = \sum_{k=1}^{d/2} \left[ (q_{2k-1} c_{2k-1} + q_{2k} c_{2k}) \cos(\Delta \omega_k) + (q_{2k} c_{2k-1} - q_{2k-1} c_{2k}) \sin(\Delta \omega_k) \right].
\end{equation}

\subsection{Refutation of the Uniform Random Rotation Assumption}

The assumption that RoPE rotates all dimensions completely randomly such that $\E[q^\top R_\Delta c] \to 0$ is mathematically false due to the decaying frequency spectrum of RoPE.

\begin{proposition}[Low-Frequency Preservation at Scale]
Let the chunk distance be $\Delta = 2048$ chunks (representing 1M tokens with 512-token chunks) for a model of width $d = 256$.
\begin{enumerate}
    \item \textbf{High-frequency subspaces} (small $k$, e.g., $k = 1$):
    \begin{equation}
    \omega_1 = 1 \implies \theta_1(2048) = 2048 \text{ radians} \approx 325.9 \text{ full rotations},
    \end{equation}
    causing extreme phase wrapping.
    \item \textbf{Low-frequency subspaces} (large $k$, e.g., $k = 128$):
    \begin{equation}
    \omega_{128} = 10000^{-1} = 0.0001 \implies \theta_{128}(2048) = 0.2048 \text{ radians} \approx 11.7^\circ.
    \end{equation}
    At $11.7^\circ$, the vectors remain highly aligned ($\cos(11.7^\circ) \approx 0.98$), refuting uniform random rotation.
\end{enumerate}
\end{proposition}

\subsection{Derivation of Frequency OOD Phase Scrambling}

The collapse of retrieval accuracy is caused by Frequency OOD Phase Scrambling of high-frequency dimensions.

\begin{theorem}[Frequency OOD Phase Scrambling]
Assume a retriever is trained on short contexts ($\Delta \leq 16$ chunks), where the maximum rotation satisfies $\theta_k(\Delta) \leq 16$ radians, allowing the model to learn Query/Key projections $W_Q, W_K$ that associate specific semantic features with small, highly coherent phase shifts. Under out-of-distribution extrapolation ($\Delta \geq 2048$), middle-to-high frequency dimensions ($k < d/4$) experience extreme phase wrapping:
\begin{equation}
\theta_k(\Delta) \in [16, 2048] \text{ radians},
\end{equation}
rotating vectors through hundreds of full cycles. The rotated semantic score in these subspaces behaves as:
\begin{equation}
S_k(i, j) \approx u_k \cdot \text{Uniform}(-1, 1),
\end{equation}
where $u_k = (q_{2k-1} c_{2k-1} + q_{2k} c_{2k})$. This completely scrambles the high-resolution keyword matching coordinate space, overwhelming the low-capacity, slowly-rotated low-frequency dimensions and causing retrieval failure.
\end{theorem}

\subsection{The Content-Position Separation Corollary}

To prevent Frequency OOD Phase Scrambling, we formulate the Content-Position Separation Corollary.

\begin{corollary}[Content-Position Separation]
For any retriever operating on sequences longer than its maximum training length ($\Delta_{\text{test}} \gg \Delta_{\text{train}}$), semantic matching scores must be computed in a purely position-invariant (content-only) space:
\begin{equation}
S(i, j) = (W_Q q)^\top (W_K c_j),
\end{equation}
with position processing deferred entirely to the generative stage. This preserves a stable signal-to-noise ratio:
\begin{equation}
\text{SNR}_{\text{content}} \approx \mathcal{O}\left(\frac{d}{\ln N}\right).
\end{equation}
\end{corollary}

\section{Optimizer Dynamics for Guard Awakening}
\label{app:optimizer}

\begin{proposition}[LayerScale attenuation of branch learning]
\label{prop:layerscale-grad}
Consider a guarded residual branch
\begin{equation}
y=x+\lambda B_\phi(x),
\end{equation}
and let $g=\partial\mathcal{L}/\partial y$. Then
\begin{equation}
\frac{\partial\mathcal{L}}{\partial\lambda}=\langle g,B_\phi(x)\rangle,
\qquad
\nabla_\phi\mathcal{L}=\lambda\,J_{B_\phi}(x)^\top g.
\end{equation}
Thus a very small $\lambda$ protects the residual stream but also attenuates learning inside $B_\phi$. The scale parameter can still receive an unattenuated gradient, but only if the untrained branch output $B_\phi(x)$ is already aligned with the downstream loss gradient.
\end{proposition}

\begin{proof}
The equations follow directly from differentiating $y=x+\lambda B_\phi(x)$ with respect to $\lambda$ and $\phi$ and applying the chain rule.
\end{proof}

Consider a guarded branch $x\mapsto x+\lambda B_\phi(x)$ with $|\lambda|\ll1$ at initialization. Proposition~\ref{prop:layerscale-grad} explains the practical unused-branch risk. If the backbone quickly reduces the loss, this alignment signal can vanish before the branch learns a useful representation.

\begin{proposition}[Effect of global norm clipping]
\label{prop:global-clipping}
Let the total gradient be partitioned into backbone and guard components, $g=(g_b,g_g)$. Global norm clipping with threshold $C$ updates both groups using the common multiplier
\begin{equation}
\alpha = \min\left(1,\frac{C}{\sqrt{\norm{g_b}_2^2+\norm{g_g}_2^2}}\right).
\end{equation}
If $\norm{g_b}_2\gg \norm{g_g}_2$ and $\norm{g_b}_2>C$, then the guard update is approximately $C g_g/\norm{g_b}_2$, even when $g_g$ is already small. Per-group clipping replaces this multiplier by $\alpha_g=\min(1,C_g/\norm{g_g}_2)$ for the guard group and removes the direct dependence on the backbone norm.
\end{proposition}

\begin{proof}
The global clipping formula follows from the definition of norm clipping. The approximation follows by neglecting $\norm{g_g}_2^2$ in the denominator when the backbone norm dominates. Per-group clipping applies the same operation to each partition separately.
\end{proof}

The remedy is a curriculum and measurement protocol rather than an optimizer myth:
\begin{enumerate}
    \item train or pretrain the branch on auxiliary retrieval/local-pattern objectives;
    \item briefly freeze or slow the backbone so the branch receives meaningful gradients;
    \item use per-group learning rates and per-group clipping;
    \item anneal $\lambda$ from a safe value to a useful value;
    \item measure branch usage by causal-drop tests and gradient norms.
\end{enumerate}

\end{document}
\typeout{get arXiv to do 4 passes: Label(s) may have changed. Rerun}
